%==============================================================================
%  Interrogation n°2 - Mathematiques - Classe de 5e - CORRIGE
%  Theme : Angles (adjacents, complementaires, supplementaires)
%  Compilation : pdflatex
%==============================================================================
\documentclass[12pt,a4paper]{article}

\usepackage[utf8]{inputenc}
\usepackage[T1]{fontenc}
\usepackage{lmodern}
\usepackage[french]{babel}
\usepackage[a4paper,margin=2cm]{geometry}
\usepackage{amsmath,amssymb}
\usepackage{array}
\usepackage{tabularx}
\usepackage[table]{xcolor}
\definecolor{entete}{RGB}{225,235,225}
\renewcommand{\arraystretch}{1.5}
\newcolumntype{C}[1]{>{\centering\arraybackslash}p{#1}}

\pagestyle{empty}
\setlength{\parindent}{0pt}

\begin{document}

%------------------------------------------------------------------ EN-TETE
\begin{center}
\begin{tabularx}{\textwidth}{|X|C{4.2cm}|}
\hline
\textbf{Lycée Mixte 1 de Yamoussoukro} & \textbf{Année scolaire : 2026-2027} \\
\textbf{Discipline : Mathématiques} & \textbf{Classe : 5\textsuperscript{e}} \\
\hline
\end{tabularx}
\end{center}

\vspace{0.3cm}
\begin{center}
{\Large\bfseries CORRIGÉ DE L'INTERROGATION N\textsuperscript{o} 2}\\[2pt]
{\large \textit{Thème : Les angles}}
\end{center}

\vspace{0.2cm}
\textbf{Rappels de cours.}
\begin{itemize}
  \item \textbf{Angles adjacents :} deux angles qui ont le même sommet, un côté commun, et qui
        sont situés de part et d'autre de ce côté commun.
  \item \textbf{Angles complémentaires :} deux angles dont la somme des mesures est égale à 90\degres.
  \item \textbf{Angles supplémentaires :} deux angles dont la somme des mesures est égale à 180\degres.
\end{itemize}

\vspace{0.3cm}

%------------------------------------------------------------------ TABLEAU CORRIGE
\begin{center}
\begin{tabularx}{\textwidth}{|C{0.7cm}|X|C{1.6cm}|X|}
\hline
\rowcolor{entete}
\textbf{N\textsuperscript{o}} & \textbf{Affirmation} & \textbf{Réponse} & \textbf{Justification} \\
\hline
1 & Deux angles adjacents ont le même sommet, un côté commun et sont situés de part et d'autre de ce côté commun. & \textbf{VRAI} &
C'est exactement la définition des angles adjacents (les trois conditions). \\
\hline
2 & Si deux angles ont seulement un sommet commun, alors ils sont adjacents. & \textbf{FAUX} &
Il faut aussi un côté commun et être de part et d'autre de ce côté. \\
\hline
3 & Deux angles adjacents sont situés de part et d'autre de leur côté commun. & \textbf{VRAI} &
C'est une condition de la définition. \\
\hline
4 & Deux angles complémentaires ont pour somme 90\degres. & \textbf{VRAI} &
C'est la définition d'angles complémentaires. \\
\hline
5 & Deux angles supplémentaires ont pour somme 90\degres. & \textbf{FAUX} &
La somme de deux angles supplémentaires est 180\degres, pas 90\degres. \\
\hline
6 & 35\degres{} et 55\degres{} sont complémentaires. & \textbf{VRAI} &
$35\degres + 55\degres = 90\degres$. \\
\hline
7 & 120\degres{} et 70\degres{} sont supplémentaires. & \textbf{FAUX} &
$120\degres + 70\degres = 190\degres \neq 180\degres$. \\
\hline
8 & Le supplément de 60\degres{} mesure 120\degres. & \textbf{VRAI} &
$180\degres - 60\degres = 120\degres$. \\
\hline
9 & Deux angles complémentaires peuvent être tous les deux obtus. & \textbf{FAUX} &
Un angle obtus dépasse 90\degres ; deux tels angles dépasseraient 90\degres à eux seuls. \\
\hline
10 & Le complément de 45\degres{} mesure 45\degres. & \textbf{VRAI} &
$90\degres - 45\degres = 45\degres$. \\
\hline
\end{tabularx}
\end{center}

\vspace{0.3cm}
\textbf{Barème :} 1 point par bonne réponse, soit un total de 10 points.

\vspace{0.2cm}
\textbf{Corrigé en bref :}
1-V \quad 2-F \quad 3-V \quad 4-V \quad 5-F \quad 6-V \quad 7-F \quad 8-V \quad 9-F \quad 10-V

\end{document}
